% Full source-specific arguments; shared drivers are integrated by the lead.

\begin{annexitem}{annex:p9.n-c-023.1.positive-closure}{thm:p09-gl88-positive-closure}{Positive closure, leastness and atomic consequence}
\begin{proof}
The operator \(T_Q\) is monotone: every finite rule body contained in
\(X\subseteq Y\) is contained in \(Y\).  Since
\(X_0=\varnothing\subseteq T_Q(X_0)=X_1\), induction gives
\(X_n\subseteq X_{n+1}\) for every \(n\).  Put
\(X_\omega=\bigcup_n X_n\).

If a rule body is contained in \(X_\omega\), its finitely many atoms all
occur by some common finite stage \(X_n\); the head belongs to
\(X_{n+1}\subseteq X_\omega\).  This includes the empty body.  Hence
\(T_Q(X_\omega)\subseteq X_\omega\).  Conversely, any member of
\(X_\omega\) belongs to some \(X_{n+1}=T_Q(X_n)\), so monotonicity puts
it in \(T_Q(X_\omega)\).  Thus \(X_\omega\) is fixed and satisfies every
rule of \(Q\).

For any model \(N\) of \(Q\), \(T_Q(N)\subseteq N\).  Starting from
\(X_0\subseteq N\), induction gives \(X_n\subseteq N\); therefore
\(X_\omega\subseteq N\).  It is the least model and consequently the unique
minimal model.  These arguments do not require \(Q\) or its atom set to be
finite; finiteness of each individual rule body is the necessary union step.

Each atom in \(X_\omega\) lies in every propositional model of the ground
implication theory \(Q\), hence is entailed.  If an atom is not in
\(X_\omega\), the valuation making precisely \(X_\omega\) true satisfies
\(Q\) and falsifies that atom, so it is not entailed.  Atoms in a larger
\(At\) which are absent from \(Q\) can likewise be assigned false.  This
proves \(th(Q)\cap At=X_\omega=M_Q\).  In particular the empty program
has least model \(\varnothing\), and a positive program has that one least
model as its only stable set.  This is the dependency invoked by the source's
\S2 and Theorem~3 \cite[pp.~4, 9]{GelfondLifschitz1988}.
\end{proof}
\end{annexitem}

\begin{annexitem}{annex:p9.n-c-023.1.theorem1}{thm:p09-gl88-theorem1}{The source's modelhood and minimality argument}
\begin{proof}
Let \(M=S_\Pi(M)\), so \(M\) is the least model of \(\Pi_M\).
Consider an original rule \(R\).  If a negative body atom \(B\) belongs
to \(M\), its literal \(\neg B\) is false in \(M\), so the original
implication is satisfied.  Otherwise every negative literal is true in
\(M\), and the positive rule \(R'\) obtained by erasing them is in
\(\Pi_M\).  Since \(M\) satisfies \(R'\), it satisfies \(R\).
Therefore \(M\) models \(\Pi\).

Suppose \(N\subseteq M\) is also a model of \(\Pi\).  For every rule
\(R'\) in \(\Pi_M\), take its original rule \(R\).  Each negative atom
of \(R\) is outside \(M\), hence outside \(N\).  Its negative literals
are therefore all true in \(N\).  If the positive body of \(R'\) is true
in \(N\), then all of \(R\)'s body is true, so its head is true because
\(N\) models \(R\).  Thus \(N\) models every rule of \(\Pi_M\).
Leastness of \(M\) gives \(M\subseteq N\), and therefore \(N=M\).
No proper subset of \(M\) models \(\Pi\).

This reconstructs both stages of original Theorem~1's supplied proof.  It
does not assert the false converse that every minimal original model is
stable \cite[\S2, Theorem~1 and proof, pp.~4--5]{GelfondLifschitz1988}.
\end{proof}
\end{annexitem}

\begin{annexitem}{annex:p9.n-c-023.1.positive-body}{thm:p09-gl88-positive-body}{No positive derivation can start}
\begin{proof}
Every surviving rule of every \(\Pi_M\) still has at least one positive
body atom.  Hence \(T_{\Pi_M}(\varnothing)=\varnothing\), and every
positive closure stage is empty.  Consequently
\(S_\Pi(M)=\varnothing\) for all \(M\); its only fixed point is
\(\varnothing\).  This includes program (6), regardless of which negative
conditions delete its rules \cite[\S3, pp.~6--7]{GelfondLifschitz1988}.
\end{proof}
\end{annexitem}

\begin{annexitem}{annex:p9.n-c-023.1.objective-expansion}{thm:p09-gl88-objective-expansion}{Constructing the entire expansion by finite modal depth}
\begin{proof}
All formulas are finite.  Modal depth is zero for objective atoms, is
unchanged by negation, takes the maximum over Boolean arguments, and
increases by one under \(L\).  Boolean constants \(\top,\bot\) below
can be read as fixed objective tautology and contradiction abbreviations.

For the objective theory \(A\), define an objective reduction \(r_A(F)\)
by structural recursion.  Leave objective atoms unchanged, commute with
every Boolean connective, and put
\[
 r_A(LG)=
 \begin{cases}
 \top,&A\models r_A(G),\\
 \bot,&A\not\models r_A(G).
 \end{cases}
\]
The test in this clause uses the already constructed reduction of the
strictly smaller formula \(G\).  It is therefore well defined; no
fixed-point existence is hidden in this definition.  Set
\[
 E_A=\{F:A\models r_A(F)\}.
\]

First suppose \(A\) is consistent.  Every ordinary objective valuation
\(v\) satisfying \(A\) has an extension to the Boolean skeleton of the
modal language: assign each propositional atom of the form \(LG\) value
one exactly when \(G\in E_A\).  Modal atoms are independent letters for
propositional consequence; no Kripke condition or modal axiom is being
imposed.  This extension satisfies precisely the required sign for every
\(LG\) in \(\Delta(E_A)\).  Conversely, every valuation satisfying
\(A\cup\Delta(E_A)\) restricts to an objective model of \(A\) and has
these forced values for every modal atom.  Structural induction yields
\[
 \widehat v(F)=v(r_A(F))
\]
for every formula \(F\).  Consequently
\[
 A\cup\Delta(E_A)\models F
 \quad\Longleftrightarrow\quad A\models r_A(F)
 \quad\Longleftrightarrow\quad F\in E_A.
\]
Thus \(E_A=th(A\cup\Delta(E_A))\), establishing existence and the full
expansion equation, not just its restriction to atoms.

For uniqueness, let \(E\) be any stable expansion of this same consistent
\(A\).  The literals \(\Delta(E)\) assign exactly one sign to each
modal atom \(LG\), so any objective model of \(A\) extends to a model
of \(A\cup\Delta(E)\).  In particular that union is consistent.  Its
objective consequences are exactly those of \(A\), proving
\(F\in E\) iff \(A\models F\) for every depth-zero formula \(F\).

Induct on modal depth.  Suppose membership in \(E\) agrees with membership
in \(E_A\) for every formula of depth less than \(d\).  In a formula
\(F\) of depth at most \(d\), each modal atom \(LG\) in its Boolean
skeleton has \(G\) of depth less than \(d\).  Its forced value under
\(\Delta(E)\) is therefore exactly the value prescribed in
\(r_A(LG)\).  The same valuation extension argument now gives
\[
 F\in E
 \ \Longleftrightarrow\ A\cup\Delta(E)\models F
 \ \Longleftrightarrow\ A\models r_A(F)
 \ \Longleftrightarrow\ F\in E_A.
\]
The induction covers every finite formula, so \(E=E_A\).

If \(A\) is inconsistent, every objective formula follows from \(A\),
and \(E_A\) is the whole modal language.  Also
\(th(A\cup\Delta(E))\) is the whole language for every \(E\), so the
whole language is the only fixed point.  This handles the boundary case
without silently assuming consistency.  Finally \(r_A(F)=F\) for every
objective \(F\), proving objective conservativity in both cases.

The source records this result as a fact cited to [8] and [6].  The
modal-depth construction here is an explicitly supplied editorial proof of
that precise fact; it is not claimed to reproduce the unavailable wording
of those earlier proofs \cite[\S4, p.~8]{GelfondLifschitz1988}.
\end{proof}
\end{annexitem}

\begin{annexitem}{annex:p9.n-c-023.1.expansion-reduct}{thm:p09-gl88-expansion-reduct}{Equivalence under all introspective assumptions}
\begin{proof}
Fix an arbitrary \(E\) and put \(M=E\cap At\).  The same assumption
set \(\Delta(E)\) is used on both sides.  Consider a translated rule
\[
 R_I=(A_1\land\cdots\land A_m\land
       \neg LB_1\land\cdots\land\neg LB_n)\supset H.
\]
If some \(B_j\in M\), then \(B_j\in E\) and \(LB_j\in\Delta(E)\).
In every valuation satisfying \(\Delta(E)\), the antecedent of \(R_I\)
is false.  The formula imposes no restriction and is equivalent to
\(\top\) under these assumptions; this is exactly the rule-deletion
case in \(\Pi_M\).

If all \(B_j\notin M\), each \(B_j\) is an atom of \(At\) not in
\(E\), so \(\neg LB_j\in\Delta(E)\).  The translated antecedent is
therefore equivalent, under these assumptions, to its positive part.
The rule is equivalent to
\((A_1\land\cdots\land A_m)\supset H\), exactly the retained
rule of \(\Pi_M\).  Empty positive or negative parts pose no problem:
the empty conjunction is \(\top\).

Rule by rule, even for infinitely many ground rules, the two theories
\[
 I(\Pi)\cup\Delta(E),\qquad \Pi_M\cup\Delta(E)
\]
have the same propositional models and hence the same closure.  Therefore
\[
 E=th(I(\Pi)\cup\Delta(E))
 \quad\Longleftrightarrow\quad
 E=th(\Pi_{E\cap At}\cup\Delta(E)).
\]
This proves both directions of the original lemma.  No presupposition that
\(E\) is already an expansion, nor any replacement of the all-formula
assumptions by only atomic ones, was needed
\cite[\S4, lemma and proof, p.~9]{GelfondLifschitz1988}.
\end{proof}
\end{annexitem}

\begin{annexitem}{annex:p9.n-c-023.1.theorem3}{thm:p09-gl88-theorem3}{Przymusinska's two-stage existence and uniqueness argument}
\begin{proof}
Let \(M\) be the unique stable model of \(\Pi\).  Its reduct \(\Pi_M\)
is objective, so Lemma~\ref{thm:p09-gl88-objective-expansion} supplies
its unique expansion \(E\).  Objective conservativity and
Lemma~\ref{thm:p09-gl88-positive-closure} give, in the source notation,
\[
 E\cap At=th(\Pi_M)\cap At=S_\Pi(M)=M.
\]
The candidate-dependent reduct is therefore
\(\Pi_{E\cap At}=\Pi_M\).  By the expansion--reduct lemma, \(E\)
is a stable expansion of \(I(\Pi)\), establishing existence.

Now let \(E'\) be any stable expansion of \(I(\Pi)\).  The lemma
makes it an expansion of the objective theory \(\Pi_{E'\cap At}\).
Again using objective conservativity and positive closure,
\[
 E'\cap At=th(\Pi_{E'\cap At})\cap At
          =S_\Pi(E'\cap At).
\]
Thus \(E'\cap At\) is a stable set of \(\Pi\).  Uniqueness of that
set gives \(E'\cap At=M\), so both \(E\) and \(E'\) are expansions
of exactly the same objective reduct \(\Pi_M\).  Its unique expansion
forces \(E'=E\).

This is the complete argument supplied for original Theorem~3 and expressly
credited there to Halina Przymusinska.  The formerly external
objective-expansion step has been proved separately above rather than
renamed as a premise \cite[\S4, pp.~8--9]{GelfondLifschitz1988}.
\end{proof}
\end{annexitem}

\begin{annexitem}{annex:p9.n-c-023.1.correspondence}{thm:p09-gl88-correspondence}{Projection and reconstruction are inverse on all stable sets}
\begin{proof}
For any stable expansion \(E\) of \(I(\Pi)\), the same lemma and
objective-conservativity calculation used in Theorem~3 give
\(E\cap At=S_\Pi(E\cap At)\).  Thus projection gives a stable set.
For any stable set \(M\), let \(E_M\) be the unique expansion of
\(\Pi_M\).  Then
\(E_M\cap At=th(\Pi_M)\cap At=S_\Pi(M)=M\); the lemma makes
\(E_M\) an expansion of \(I(\Pi)\).

Projection followed by reconstruction returns the starting expansion,
because any expansion projecting to \(M\) is an expansion of
\(\Pi_M\), whose expansion is unique.  Reconstruction followed by
projection returns \(M\) by the displayed equality.  Hence the two maps
are inverse, including the case in which there are no stable sets.
Zero stable models give zero expansions, and the two-model source example
gives exactly two expansions.  This general correspondence follows from
the supplied arguments but is separated from the original theorem's
explicit unique-model hypothesis \cite[\S4, pp.~8--9]{GelfondLifschitz1988}.
\end{proof}
\end{annexitem}
